\RequirePackage{luatex85}
\documentclass[leqno]{ltjsarticle}
\usepackage{luatexja-fontspec}
\usepackage[top=20truemm,bottom=20truemm,left=20truemm,right=20truemm]{geometry}
\usepackage{luatexja} 
\usepackage{multicol,amsmath,amssymb,mathtools,ascmac,amsthm,amscd,physics,comment,dcolumn,titlesec,mathrsfs,mypkg}
\usepackage[all]{xy}
\titleformat*{\section}{\Large\bfseries}
\setlength{\parindent}{0pt}
\pagestyle{plain}
%\everymath{\displaystyle}
\setcounter{page}{3}
\begin{document}
\hfill (PD申請内容ファイル)\\
\fbox{\parbox{\textwidth}{
		{\large【2】研究計画}\ {\small 適宜概念図を用いるなどして、わかりやすく記入してください。様式の変更・追加は不可です。 }
}}\\ \\
\fbox{\parbox{\textwidth}{
		(1) 研究の概要及び研究の位置づけ 本項目は 1 頁に収めてください。\\
・まず、研究課題名及び研究の概要を全角 500 字（半角 1,000 字）程度で記入してください。\\
・続けて、特別研究員として取り組む研究の位置づけについて、当該分野の状況や課題等の背景、並びに本研究計画の着想に至った経緯も含めて記入してください。
}}\\ \\
\underline{研究課題名 ： 軌道体射影直線の導来圏の半直交分解とBridgeland安定性条件に関する研究}
\\
\bullet\  研究の概要\\
本研究では，軌道体射影直線の導来圏に対する Bridgeland 安定性条件の空間を具体的に記述することを目的とする．Bridgeland 安定性条件は，導来圏の対象を「安定な構成要素」に分解する理論であり，その全体は複素多様体を成し，自己同値群の作用，モジュライ空間の壁越え，ホモロジー的ミラー対称性と深く関わる重要な幾何学的対象である．

軌道体射影直線の導来圏には半直交分解が存在し，その成分として ADE 型箙の導来圏が現れる．ADE 型箙の導来圏については，安定性条件空間と配置空間や組紐群との関係が詳しく研究されている．そこで本研究では，各成分の部分圏上で既に理解されている安定性条件から，全体の導来圏上の安定性条件を構成するための条件を定式化する．さらに，この方法を用いて，安定性条件空間の連結性，壁と部屋の構造，自己同値群の作用を記述する．

本研究により，半直交分解という離散的構造から安定性条件空間という連続的な幾何学的構造を復元する方法を確立し，表現論，代数幾何，ミラー対称性を結びつける新たな枠組みを与えることを目指す．

\vspace{2mm}

\bullet\ 研究の位置づけ\\
Bridgeland 安定性条件は，導来圏の対象を「安定な構成要素」に分解する理論であり，
その全体は複素多様体を成す．この空間の構造を理解することは，
自己同値群の作用，モジュライ空間の壁越え，ホモロジー的ミラー対称性を理解する上で重要である．

ADE 型箙の導来圏については，Bessis による組紐群と配置空間の研究，
Hertling による Frobenius 多様体の研究，
および Haiden--Katzarkov--Kontsevich による有限型・アフィン A 型の場合の研究により，
安定性条件空間の位相的・幾何学的構造が詳しく調べられている．
これらの研究では，安定性条件空間が配置空間の被覆として記述され，
組紐群や Coxeter 変換との深い関係が明らかにされている．

一方，軌道体射影直線の導来圏については，
その内部に ADE 型箙の導来圏が部分圏として現れるにもかかわらず，
部分圏上の安定性条件から全体の安定性条件をどのように構成するかは
未解決である．
すなわち，半直交分解という導来圏の離散的構造から，
安定性条件空間という連続的な幾何学的構造をどこまで復元できるかは，
導来圏論における基本的な問題である．

本研究は，軌道体射影直線を具体例として，
半直交分解に沿った安定性条件の貼り合わせ条件を定式化し，
安定性条件空間の連結性，壁と部屋の構造，自己同値群の作用を記述する．
これにより，ADE 型で確立された理論を
より一般の Deligne--Mumford スタックへ拡張し，
表現論，代数幾何，ミラー対称性を結びつける新たな枠組みを与える．
\newpage
\fbox{\parbox{\textwidth}{
	{\large【2】研究計画\ (続き)}\ {\small 適宜概念図を用いるなどして、わかりやすく記入してください。様式の変更・追加は不可です。 }
}}\\ \\
\fbox{\parbox{\textwidth}{
		(2) 研究目的・内容等 本項目は２頁に収めてください。
① 特別研究員として取り組む研究計画における研究目的、研究方法、研究内容について記入してください。
② どのような計画で、何を、どこまで明らかにしようとするのか、特別研究員奨励費の応募区分（下記（※）参照）に応じて、年次計画を示
し、具体的に記入してください。研究計画が想定通り進まなかった場合の対応方法があれば、あわせて記入してください。
③ 研究の特色・独創的な点にも触れて記入してください。（研究の特色・独創的な点の例：先行研究等との比較、本研究の完成時に予想され
るインパクト、総合知への貢献、将来の見通し）
④ 研究計画が受入研究室としての研究活動や共同研究の一部と位置づけられる場合は申請者が担当する部分を明らかにしてください。
⑤ 研究計画の期間中に受入研究機関と異なる研究機関（外国の研究機関等を含む。）において研究に従事することも計画している場合は、具
体的に記入してください。
（※）特別研究員奨励費の研究期間が 3 年の場合の応募総額は（A 区分）が 300 万円以下、（B 区分）が 300 万円超 450 万円以下。2 年の場合は（A 区分）が 200 万円以下、（B
区分）が 200 万円超 300 万円以下。1 年の場合は（A 区分）が 100 万円以下、（B 区分）が 100 万円超 150 万円以下。（B 区分については研究計画上必要な場合のみ記入）
}}\\ \\
本研究の目的は，軌道体射影直線の導来圏に対する Bridgeland 安定性条件の空間を具体的に記述することである．
特に，導来圏に現れる半直交分解に着目し，各成分の部分圏上で既に理解されている安定性条件から，
全体の導来圏上の安定性条件を構成するための条件を与えることを目指す．
これにより，安定性条件空間の連結性，壁と部屋の構造，自己同値群の作用を記述する．

軌道体射影直線の導来圏には例外対象列や tilting 対象が存在し，
その部分圏として ADE 型箙の導来圏が現れる．
ADE 型箙の導来圏については，
安定性条件空間と配置空間や組紐群との関係が詳しく研究されている．
本研究では，これらの既知の結果を出発点とし，
軌道体射影直線の導来圏における半直交分解の構造を明示的に記述する．

次に，各成分の部分圏上の安定性条件を比較し，
それらが全体の導来圏上で Harder--Narasimhan 分解を与えるための
整合条件を定式化する．
特に，部分圏をまたぐ拡大構造と，
例外対象列の変異によって生じる半直交分解の変化を解析し，
安定性条件を貼り合わせる一般的な方法を構築する．

さらに，得られた方法を用いて，
軌道体射影直線の安定性条件空間の連結成分，
壁と部屋の配置，自己同値群の作用を具体的に記述する．
また，Bessis や Hertling の研究，
および Haiden--Katzarkov--Kontsevich による結果と比較することで，
安定性条件空間と配置空間，Frobenius 多様体，
ミラー対称性との対応を検証する．

以上により，半直交分解という導来圏の離散的構造から，
安定性条件空間という連続的な幾何学的構造を復元する方法を確立し，
表現論，代数幾何，ミラー対称性を結びつける新たな枠組みを与える．
\newpage
（【２】研究計画（２）研究目的・内容等の続き）
\newpage
\fbox{\parbox{\textwidth}{
	{\large【2】研究計画\ (続き)}\ {\small 適宜概念図を用いるなどして、わかりやすく記入してください。様式の変更・追加は不可です。 }
}}\\ \\
\fbox{\parbox{\textwidth}{
		(3) 受入研究室の選定理由 本項目は１頁に収めてください。
 採用後の受入研究室を選定した理由について、次の項目を含めて記入してください。
 ① 受入研究室を知ることとなったきっかけ、及び、採用後の研究実施についての打合せ状況
 ② 申請の研究課題を遂行するうえで、当該受入研究室で研究することのメリット、新たな発展・展開
 ※ 個人的に行う研究で、指導的研究者を中心とするグループが想定されない分野では、「研究室」を「研究者」と読み替えて記入して
ください。
}}\\ \\
\newpage
\fbox{\parbox{\textwidth}{
		【３】人権の保護及び法令等の遵守への対応 本項目は１頁に収めてください。様式の変更・追加は不可です。
・本欄には、「【２】研究計画」を遂行するにあたって、相手方の同意・協力を必要とする研究、個人情報の取り扱いの配慮を必要とする研究、
生命倫理・安全対策に対する取組を必要とする研究や安全保障貿易管理を必要とする研究など指針・法令等（国際共同研究を行う国・地域
の指針・法令等を含む）に基づく手続が必要な研究が含まれている場合、講じる対策と措置を記入してください。
・例えば、個人情報を伴うアンケート調査・インタビュー調査・行動調査（個人履歴・映像を含む）、国内外の文化遺産の調査等、提供を受け
た試料の使用、侵襲性を伴う研究、インフォームド・コンセントが必要な研究、ヒト遺伝子解析研究、遺伝子組換え実験、動物実験、機微技
術に関わる研究など、研究機関内外の情報委員会や倫理委員会等における承認手続が必要となる調査・研究・実験などが対象となりますの
で手続の状況も具体的に記入してください。
・なお、該当しない場合には、その旨記入してください。
}}\\ \\
特に該当なし．
\newpage
\fbox{\parbox{\textwidth}{
		【４】研究遂行力の自己分析 本項目は２頁に収めてください。様式の変更・追加は不可です。
・日本学術振興会特別研究員制度は、我が国の学術研究の将来を担う創造性に富んだ研究者の養成・確保に資することを目的としています。
この目的に鑑み、これまで携わった研究活動における経験などを踏まえ、研究遂行力について分析してください。 
}}\\ \\
\newpage
（【４】研究遂行力の自己分析の続き）
\end{document}

